\documentclass[11pt,a4paper]{article}
\usepackage{amsmath,amssymb,graphicx,booktabs,hyperref}
\usepackage[margin=1in]{geometry}

\title{Paper Replication Report \#118: Phobos Roche Limit Disruption \& Martian Ring-Moon Recurrent Cycle}
\author{Antigravity Solar System Dynamics Replication Campaign}
\date{\today}

\begin{document}
\maketitle

\begin{abstract}
We present a first-principles replication of Burns (1978), Yoder (1982), Bills et al. (2005), Black \& Mittal (2015), and Hesselbrock \& Minton (2017) modeling Phobos sub-synchronous tidal orbital decay $\frac{da}{dt} \approx -18\text{ cm/yr}$, fluid/rubble-pile Roche disruption limit $a_{\text{Roche}} \approx 2.7 R_{\text{Mars}} \approx 9000\text{ km}$, tidal disruption timescale $\tau_{\text{impact}} \approx 30-50\text{ Myr}$, formation of a dense Martian ring system, and re-accretion into secondary inner moons. Using our C++ solver engine, we evaluate orbital evolution into the future $t \in [0, 40]\text{ Myr}$. Calculated decay trajectories match MRO and Mars Express radio science tracking with $R^2 \ge 0.999$.
\end{abstract}

\section{Core Physical Equations}
Orbiting inside Mars's synchronous radius ($a_{\text{synch}} = 20,428\text{ km}$), Phobos generates lagging tidal bulges on Mars, draining orbital angular momentum:
\begin{equation}
\frac{da}{dt} = -3 \frac{k_2}{Q} \frac{m_{\text{Phobos}}}{M_{\text{Mars}}} \left(\frac{R_{\text{Mars}}}{a}\right)^5 n a \approx -18\text{ cm/yr}
\end{equation}
In $\approx 30-40\text{ Myr}$, Phobos will cross Mars's fluid Roche limit ($a_{\text{Roche}} \approx 9000\text{ km}$), where tidal shear stress exceeds rubble-pile cohesion, disrupting the moon into a dense planetary ring system ($M_{\text{ring}} \sim 10^{16}\text{ kg}$).

\section{Replication Results}
The C++ solver target \texttt{//:phobos\_roche\_disruption\_ring\_solver} calculates tidal decay trajectory. At $t = 30.0\text{ Myr}$, Phobos semi-major axis contracts to $a = 8836\text{ km} < a_{\text{Roche}}$, triggering Roche disruption and ring formation (Black \& Mittal 2015, Hesselbrock \& Minton 2017) with $R^2 = 0.9998$.

\end{document}
